
\documentclass[aps,preprint]{revtex4}
%%%%%%%%%%%%%%%%%%%%%%%%%%%%%%%%%%%%%%%%%%%%%%%%%%%%%%%%%%%%%%%%%%%%%%%%%%%%%%%%%%%%%%%%%%%%%%%%%%%%%%%%%%%%%%%%%%%%%%%%%%%%
%TCIDATA{OutputFilter=LATEX.DLL}
%TCIDATA{Created=Saturday, April 13, 2002 16:32:44}
%TCIDATA{LastRevised=Sunday, April 14, 2002 22:52:04}
%TCIDATA{<META NAME="GraphicsSave" CONTENT="32">}
%TCIDATA{<META NAME="DocumentShell" CONTENT="Articles\SW\REVTeX 4">}
%TCIDATA{CSTFile=revtex4.cst}

\newtheorem{theorem}{Theorem}
\newtheorem{acknowledgement}[theorem]{Acknowledgement}
\newtheorem{algorithm}[theorem]{Algorithm}
\newtheorem{axiom}[theorem]{Axiom}
\newtheorem{claim}[theorem]{Claim}
\newtheorem{conclusion}[theorem]{Conclusion}
\newtheorem{condition}[theorem]{Condition}
\newtheorem{conjecture}[theorem]{Conjecture}
\newtheorem{corollary}[theorem]{Corollary}
\newtheorem{criterion}[theorem]{Criterion}
\newtheorem{definition}[theorem]{Definition}
\newtheorem{example}[theorem]{Example}
\newtheorem{exercise}[theorem]{Exercise}
\newtheorem{lemma}[theorem]{Lemma}
\newtheorem{notation}[theorem]{Notation}
\newtheorem{problem}[theorem]{Problem}
\newtheorem{proposition}[theorem]{Proposition}
\newtheorem{remark}[theorem]{Remark}
\newtheorem{solution}[theorem]{Solution}
\newtheorem{summary}[theorem]{Summary}
\newenvironment{proof}[1][Proof]{\noindent\textbf{#1.} }{\ \rule{0.5em}{0.5em}}
\input{tcilatex}

\begin{document}

\preprint{HEP/123-qed}
\title[Short title for running header]{Long title that appears on the title
page}
\author{First Author}
\affiliation{My Institution}
\author{Second Author}
\affiliation{My Institution}
\author{Third Author}
\affiliation{Other Institution}
\keywords{one two three}
\pacs{PACS number}

\begin{abstract}
Shell document for REV\TeX{} 4.
\end{abstract}

\volumeyear{year}
\volumenumber{number}
\issuenumber{number}
\eid{identifier}
\date[Date text]{date}
\received[Received text]{date}
\revised[Revised text]{date}
\accepted[Accepted text]{date}
\published[Published text]{date}
\startpage{101}
\endpage{102}
\maketitle
\tableofcontents

\section{CI Coefficients of AGP States}

An AGP\ state can be expanded in an unordered linearly dependent bases as 
\begin{equation}
\frac{1}{2^{N}}\sum_{1\leq j_{1},j_{2},\cdots ,j_{2N-1},j_{2N}\leq r}\gamma
_{j_{1}j_{2}}\cdots \gamma _{j_{2N-1}j_{2N}}\left| \varphi _{j_{1}}\varphi
_{j_{2}}\cdots \varphi _{j_{2N-1}}\varphi _{j_{2N}}\right\rangle
\end{equation}%
Using the symmetry property 
\begin{equation}
\left| \varphi _{j_{1}}\varphi _{j_{2}}\cdots \varphi _{j_{2N-1}}\varphi
_{j_{2N}}\right\rangle =\left( -1\right) ^{\pi \left( \sigma \right) }\left|
\varphi _{j_{\sigma \left( 1\right) }}\varphi _{j_{\sigma \left( 2\right)
}}\cdots \varphi _{j_{\sigma \left( 2N-1\right) }}\varphi _{j_{\sigma \left(
2N\right) }}\right\rangle \qquad \forall \sigma \in S_{2N}
\end{equation}%
the coefficients of configurations that can be rearranged into each other
can be collected together as the coefficient of a single ordered
configuration $\left| \varphi _{j_{1}}\cdots \varphi _{j_{2N}}\right\rangle $
where $j_{1}<j_{2}<\cdots <j_{2N-1}<j_{2N}$ as%
\begin{eqnarray}
&&\frac{1}{2^{N}}\sum_{\sigma \in S_{2N}}\gamma _{j_{\sigma \left( 1\right)
}j_{\sigma \left( 2\right) }}\cdots \gamma _{j_{\sigma \left( 2N-1\right)
}j_{\sigma \left( 2N\right) }}\left| \varphi _{j_{\sigma \left( 1\right)
}}\varphi _{j_{\sigma \left( 2\right) }}\cdots \varphi _{j_{\sigma \left(
2N-1\right) }}\varphi _{j_{\sigma \left( 2N\right) }}\right\rangle  \nonumber
\\
&=&\frac{1}{2^{N}}\sum_{\sigma \in S_{2N}}\left( -1\right) ^{\pi \left(
\sigma \right) }\gamma _{j_{\sigma \left( 1\right) }j_{\sigma \left(
2\right) }}\cdots \gamma _{j_{\sigma \left( 2N-1\right) }j_{\sigma \left(
2N\right) }}\left| \varphi _{j_{1}}\varphi _{j_{2}}\cdots \varphi
_{j_{2N-1}}\varphi _{j_{2N}}\right\rangle  \nonumber \\
&=&Pf\left( \gamma :j_{1}\cdots j_{2N}\right) \left| \varphi _{j_{1}}\varphi
_{j_{2}}\cdots \varphi _{j_{2N-1}}\varphi _{j_{2N}}\right\rangle
\end{eqnarray}%
where 
\begin{equation}
Pf\left( \gamma :j_{1}\cdots j_{2N}\right) =Pf\left( 
\begin{array}{ccccc}
0 & \gamma _{j_{1}j_{2}} & \gamma _{j_{1}j_{4}} & \cdots & \gamma
_{j_{1}j_{2N}} \\ 
\gamma _{j_{2}j_{1}} & \ddots & \gamma _{j_{3}j_{4}} & \vdots & \gamma
_{j_{3}j_{2N}} \\ 
\vdots & \vdots & \ddots & \vdots & \vdots \\ 
\vdots & \vdots & \vdots & \ddots & \gamma _{j_{2N-1}j_{2N}} \\ 
\gamma _{j_{2N}j_{1}} & \cdots & \cdots & \gamma _{j_{2N}j_{2N-1}} & 0%
\end{array}%
\right)
\end{equation}%
The symbol stands for Pfaffian, which is a function that associates a scalar
with an antisymmetric matrices. For further details on Pfaffians see

\end{document}
